\documentclass[12pt]{article}

\usepackage{ucs}     % unicode 
\usepackage[utf8x]{inputenc}  % utf-8
\usepackage[ngerman]{babel}  % new german spelling 
\usepackage{graphicx}   % use graphics 
\usepackage{fancyhdr}   % header and footer
\usepackage{setspace}
\usepackage{url}
\usepackage{tikz}
\usepackage{tikz-qtree}
\usepackage[printonlyused]{acronym}
\usepackage{a4wide}
\usepackage{natbib}	% bibstyle 
\usepackage[section]{placeins}	% \FloatBarrier


\setcounter{tocdepth}{2}
\begin{document}

\thispagestyle{empty} 

\begin{center}
	\topskip0pt
	\vspace*{\fill}
	
	\Huge{\textbf{Seminararbeit}}\\
	\vspace{1.5cm}
	
	\Large{\textbf{Automatische Bestimmung semantischer Rollen im deutschen
	Korpus SALSA 2.0}}\\
	\vspace{1cm}
		
	\includegraphics[scale=0.3]{images/logo_hu.png}
	\vspace{1cm}

	\begin{Large}
		Institut für Informatik \& Institut für Linguistik\\
		Humboldt-Universität zu Berlin\\
		\vspace{1.5cm}
		Robert Bärhold \& Arne Binder \\
		31. März 2014 \\ %TODO
		\vspace{1cm}
		
		\begin{table}[h]
			\Large
			\centering
			\begin{tabular}{l l}
				Seminar: & Computergestützte Analyse von Sprache\\
				Seminarleiter: & Prof. Dr. Ulf Leser\\
				 		    & Prof. Dr. Anke Lüdeling \\
				Semester: & Wintersemester 2013/14 				 	
			\end{tabular}
		\end{table}	
	\end{Large}
	\vspace*{\fill}
\end{center}


\pagenumbering{roman}
\pagestyle{fancy} %eigener Seitenstil
\fancyhf{} %alle Kopf- und Fußzeilenfelder bereinigen
\renewcommand{\headrulewidth}{0pt} %obere Trennlinie
\renewcommand{\footrulewidth}{0pt} %untere Trennlinie 
\fancyfoot[C]{\thepage} %Seitennummer

% \newpage
% \tableofcontents
% \vspace{1cm}
% \listoffigures
% \vspace{1cm}
% \listoftables

\newpage
\pagenumbering{arabic}

\section{Einleitung}
- warum SRL? Motivation
- was gibt es schon? related work
\subsection{Semantische Rollen}

Prädikate zeichnen sich dadurch aus, dass sie die Komplementationsstruktur einer
überge-ordneten syntaktischen Einheit (nämlich des Satzes) bestimmen. Zum
Beispiel verlangt ein ditransitives Verb wie \glqq{}geben\grqq{} drei weitere Satzkonstituenten mit konkreten syntaktischen Eigenschaften:

\begin{center}
	\begin{tikzpicture}
		\Tree [.geben
				[.Subjekt Paul ]
				[.{} gibt ]
				[.Akk-Objekt {die Jacke} ]
				[.Dativ-Objekt {seiner Freundin} ]
			]
			
		%\Tree [.geben
		%		[.\textbf{Subjekt} Paul ]
		%		[.{} gibt ]
		%		[.\textbf{Akk-Objekt} {die Jacke} ]
		%		[.\textbf{Dativ-Objekt} {seiner Freundin} ]
		%	]
	\end{tikzpicture}
\end{center}

Aber das Bedingtheitsverhältnis zwischen dem Verb und seinen Komplementen 
endet nicht in der Syntax. Diese sind durch das Prädikat auch hinsichtlich ihrer 
\textit{semantischen} Integrierung in den Satz determiniert. Dem Komplement \glqq{}Paul\grqq{}
wird einmal die syntaktische Funktion des Subjekts zugewiesen, es wird aber auch
durch die konkrete Semantik der Handlung \glqq{}geben\grqq{} als 
\glqq{}Geber\grqq{} charakterisiert. Man spricht von den Argumenten
eines Verbs und sagt, dass sie \textit{semantische Rollen} realisieren. So realisiert \glqq{}Paul\grqq{} im obigen Beispiel die semantische Rolle \glqq{}Geber\grqq{}.

\subsection{Frame Semantik}

Die traditionelle Semantik [TODO: evtl. welche traditionelle Sem.? Referenz?] war versucht, eine möglichst abstrakte und
allgemeine, intersprachliche Formulierung von semantischen Rollen zu
finden - bei der man z. B. nicht von \glqq{}Geber\grqq{} sondern allgemein
von \glqq{}Agens\grqq{} sprechen würde. Demgegenüber steht die von C.J. Fillmore
ab den 70er Jahren entwickelte Theorie der Frame-Semantik (FS)
[\cite{fillmore1985}]. Das Fundament der Theorie ist mit dem traditionellen
Ansatz in höchstem Maße inkompatibel, für unser Vorhaben reicht es aber, die
partikuläre Auffassung von semantischen Rollen zu erwähnen, die der
Frame-semantische Ansatz mit sich bringt.

Rollen bzw. Frame-Elemente (FE) in Termini der FS werden nicht in Bezug auf die
Argumentstruktur von Verben formuliert, sondern hinsichtlich der globalen
Situation (oder des \textit{Frames}), die durch das Prädikat evoziert wird. Aufgrund
des situationellen Charakters der Frames sind nicht nur Verben sondern auch
Nomina (\textit{Mord}, \textit{Auge}\footnote{beispielsweise als Indikator für ein Ereignis der Wahrnehmung: \glqq{}Vor seinen Augen zerrissen sie die Bücher.\grqq{}}, \textit{Frau} von jemandem) oder Adjektive (\textit{stolz} auf etwas sein) potentielle
Frame-einführende Elemente, und in diesem Sinne auch Prädikate. Nachfolgend ist der Frame \glqq{}Kritik1-salsa\grqq{} dargestellt, so wie er im framesemantisch annotierten SALSA 2.0 Korpus [\cite{rehbein_adding_2012}] zu finden ist:
\begin{quote}
A Reviewer offers his or her assessment of a work of art or performance (e.g. a play, a novel, a piece of music composition etc). The Reviewer typically bases his judgment on their own expertise and uses some set of criteria for evaluating works as to their relative merit.

\textbf{Example Sentences:}
\begin{enumerate}
\item Ich hab absichtlich auch noch keine anderen \textbf{Kritiken} gelesen und möglicherweise hab ich ein paar Dinge in dem Film übersehen.

\item Ich bin auf das Buch durch die \lbrack lobende\rbrack \textsuperscript{Valence} \textbf{Kritik} \lbrack in der FAZ\rbrack \textsuperscript{Medium} aufmerksam geworden.

\item \lbrack Olaf Storbecks Jahrhundertkrise\rbrack \textsuperscript{Evaluee} \lbrack erntet\rbrack \textsuperscript{Support} erste \lbrack gute\rbrack \textsuperscript{Valence} \textbf{Kritiken}.

\item Star Dreck - \lbrack Negative\rbrack \textsuperscript{Valence} \textbf{Kritiken} \lbrack von Fans\rbrack \textsuperscript{Reviewer}  \lbrack zum neuen Star Trek\rbrack \textsuperscript{Evaluee}
\end{enumerate}


\textbf{FEs}


\textbf{FE1(Reviewer)}: The person assessing the Evaluee for its merit.

\textbf{FE2(Evaluee)}: The cultural artifact that is evaluated.

\textbf{FE3(Medium)}: The venue or publication in which the review is distributed.

\textbf{FE4(Valence)}: The positive or negative nature of the evaluation, if specified. 
\end{quote}

\subsection{Semantic Role Labeling}

Mit Semantic Role Labeling (SRL) ist die automatisierte Erkennung und
Annotierung von semantischen Rollen innerhalb eines Satzes gemeint. Die primäre
Aufgabe von SRL ist die genaue Identifizierung der semantischen Beziehung
zwischen einem Prädikat und seinen assoziierten Elementen und Eigenschaften
(siehe \cite{SRL2008} für eine allgemeine Darstellung). In Bezug auf eine FS lässt
sich folgendes feststellen. Da FE Frame-spezifisch sind, stellen sie eine
Zwischenstufe in der semantischen Abstraktion zwischen der rein lexikalischen
Bedeutung und den traditionellen Verbenübergreifenden [@Enrique: was meinst hier mit Verb..übergr.?] semantischen Rollen dar.
Die dadurch erreichte \glqq{}mildere\grqq{} Stufe der Abstraktion über die
Rollen der verschiedenen Verben, Nomina und Adjektive eignet sich besonders gut
für die Lösung unterschiedlicher Aufgabenstellungen in Bereichen des
Sprachverstehens [TODO: ?] wie der Informationsextraktion oder Dialogsystemen [TODO: Verweis/Erklärung].

[TODO: The relationship between such surface manifestations and
semantic roles is the subject of linking theory—see \cite{levinrappaport}
for a synthesis of work in this area. In general, linking theory argues that the syntactic
realization of arguments of a predicate is predictable from semantics—exactly how this
relationship works is the subject of much debate. Regardless of the underlying mechanisms
used to generate syntax from semantics, the relationship between the two suggests
that it may be possible to learn to recognize semantic relationships from syntactic
cues, given examples with both types of information.]

Die Aufgabenstellung eines FS-orientierten SRL-Systems lässt sich 
wie folgt untergliedern [\cite{gildea}]:
\begin{enumerate}
\item Zunächst wird der Frame, der durch den Satz oder die Phrase realisiert wird, bestimmt.
\item Es folgt die Bestimmung der Frame-Elemente bezüglich des Frames. 
\item Schließlich erfolgt die Klassifizierung der einzelnen Frame-Elemente.
\end{enumerate}
\section{Problemstellung}
Im Zuge dieser Arbeit soll die Frage geklärt werden, ob sich der Ansatz von \cite{gildea} für das Deutsche adaptieren lässt(TODO: und der SALSA 2.0 Korpus eine geeignete Datenbasis darstellt?).

\subsection{Hypothese}
Die an Hand einer Baumbank, das heißt einem Phrasen-annotierten Textkorpus, der deutschen Sprache zur Verfügung stehenden lexikalischen und syntaktischen Informationen ermöglichen verlässliche Annahmen über die Realisierung von semantischen Rollen durch Konstituenten eines Satzes. Das lässt sich dazu nutzen, ein automatisiertes statistisches Lernverfahren basierend auf einem zusätzlich Rollen-annotierten Korpus zu entwickeln, welches die automatische Bestimmung der in einem deutschen Satz realisierten semantischen Rollen mit akzeptabler Genauigkeit(TODO: oder anderer Ausdruck?) ermöglicht.

\subsection{Korpus: Salsa 2.0}
Die SALSA Korpora sind für akademische Zwecke frei zugängliche, framesemantisch annotierte Korpora der deutschen Sprache. Sie basieren auf TIGER [\cite{tiger}], einer Baumbank aus deutschen Zeitungsartikeln, welche semi-automatisch mit POS-Tags und syntaktischen Strukturen annotiert wurde.\footnote{siehe http://www.ims.uni-stuttgart.de/forschung/ressourcen/korpora/tiger.html}

SALSA erweitert TIGER um semantische Informationen. Jeder Satz kann einen oder mehrere Frames besitzen. Diese bestehen aus einem Target, das ist das Prädikat, das den Frame evoziert, und verschiedenen Frame-Elementen. Das sind die Elemente des Satzes, welche eine semantische Rolle innerhalb des Frames erfüllen. Frame und Frame-Elemente werden durch einen Namen und eine eindeutige ID charakterisiert, dem Target ist lediglich das zugehörige Lemma zugewiesen. Sowohl das Target als auch die Frame-Elemente können aus mehreren Konstituenten bestehen.

Das vorige Release des SALSA Korpus, SALSA 1.0 [\cite{burchardt_salsa_2006}], enthält Frames dessen Target ein Verb ist. In SALSA 2.0 [\cite{rehbein_adding_2012}] sind außerdem durch Nomen realisierte Prädikate und deren zugehörige Frames annotiert. 

SALSA 2.0 besteht aus circa 24.000 (TODO: vs. 36251??) Sätzen. Es enthält rund 20.000 verbale und mehr als 17.000 nominale Instanzen, welche durch 648 verschiedene Targets ausgelöst werden.
\section{Implementierung des SRL-Klassifikators}
hier
\subsection{Einschränkungen}
- keine Frames über mehrere Sätze
\subsection{Features}
\subsection{Arbeitsweise des Klassifikators}
\subsubsection{Erstellung des Modells}
\subsubsection{Klassifikation der Konstituenten}
\section{Evaluation des xxxxx} %TODO

Der entwickelte Ansatz der automatischen Rollenannotation wurde nach Abschluss
der Entwickelung einer Evaluation unterzogen. Diese sollte Aufschluss über die
Qualität der entwickelten Anwendung geben, so dass ein Vergleich mit ähnlichen
Systemen gezogen werden kann.

\newpage
\bibliography{biblio}
\bibliographystyle{natdin}


\end{document}